\documentclass[11pt]{article}
\usepackage[margin=1in]{geometry}
\usepackage{amsmath,amssymb,amsfonts}
\usepackage{algorithm}
\usepackage{algpseudocode}
\usepackage[T1]{fontenc}
\usepackage[utf8]{inputenc}
\title{Algorithmic Pseudocode Sample 47}
\author{Source-backed Image2Struct algorithm sample}
\date{}
\begin{document}
\maketitle
\section{Algorithm}
This sample contains algorithmic pseudocode extracted from a source-backed LaTeX benchmark dataset. It is wrapped in a minimal article document for pdfLaTeX validation.

\begin{algorithm}[htbp]
\caption{Source-backed algorithmic procedure}
\begin{algorithm}
[H]
\caption{Square system method}
\begin{algorithmic}[1]
\State Input:\begin{itemize}
\item Algebraic system of difference equations named $\Sigma'$
\item Time measured data allowing prolongation of the system.
\item For $\bar{\mu}=\mu_1,\ldots,\mu_n$ the finite set of parameters, the data $R_{\mu_i}$ of permissible intervals for each parameter value.
\end{itemize}
\State Output: Parameter values of $\Sigma'$.
\Procedure{Detect solvability}{}
\State Define $\Sigma prolong$ as an indefinite time prolongation of $\Sigma'$.
\State Redefine $\Sigma':=[]$, an `empty system' of no equations.
\State Denote $e_1,\ldots$ the equations of $\Sigma prolong$.
\State Define $J(t)=\left[\dfrac{\partial e_i}{\partial\mu_j}\right]$ where $e_i\in \Sigma',i=1,\ldots,r$, or 0 if $\Sigma'$ is empty.
\State $i:=1$
\For{$rank(J(t)):=s<n$ i.e. is not of full rank, }
\Procedure{Move equation}{}
\State $\Sigma':=\Sigma'\cup\{e_i\}$
\State $\Sigma prolong:=\Sigma prolong\backslash\{e_i\}$
\State Compute $rank(J(t))$ (note $\Sigma'$ has updated)
\If{$rank(J(t))<s+1$} $\Sigma':=\Sigma'\backslash\{e_i\}$
\EndIf
\State $i=i+1$
\If{$rank(J(t)<n$} repeat procedure Move Equation.
\EndIf
\EndProcedure
\EndFor
\Procedure{Blackbox solver and filter}{} Note Detect Solvability runs until $J(t)$ \phantom{---------}has full rank and outputs a polynomial system $\Sigma'$ that has solutions. Run any \phantom{---------}algebraic solver.
\State Filter solutions by intersecting solution set with $R_{\mu_i}$.
\EndProcedure
\EndProcedure
\end{algorithmic}
\end{algorithm}
\end{algorithm}
\end{document}
